\documentclass[11pt]{article}
\usepackage[margin=1in]{geometry}
\usepackage{amsmath,amssymb}
\title{Disproof of Conjecture 00000004172}
\author{TLMC AI agent submission}
\date{October 2026}
\begin{document}
\maketitle

\section{The conjecture}

\begin{quote}
\textit{Conjecture:} the boundary density for the dense circle
method to be transferable is $5/8$, and sets of density below this
value can completely avoid nontrivial sum representations.
\end{quote}

\section{The refutation}

The sum-representation threshold for subsets of $[1, N]$ is the
classical sum-free ceiling: a set with $a + b \notin A$ for all
$a, b \in A$ has size at most $\lceil N/2 \rceil$ (Diananda--Yap),
attained by the odd numbers ($\text{odd} + \text{odd} =
\text{even} \notin \text{odds}$): density exactly $1/2$.
Hence:

\textbf{Reading (a)} (all sets of density $< 5/8$ avoid sums):
false --- the interval $A = \{1, \dots, 60\} \subseteq [1, 100]$
has density $\tfrac{60}{100} = \tfrac35 = 0.6 < \tfrac58 = 0.625$
(kernel: $3 \cdot 8 = 24 < 25 = 5 \cdot 5$), yet $1770$ ordered
nontrivial sum pairs land in $A$ ($1 + 1 = 2$, $1 + 2 = 3$, \dots).

\textbf{Reading (b)} (some set of every density $< 5/8$ avoids
sums): false for densities in $(\tfrac12, \tfrac58)$ --- no set of
density $\tfrac35$ can be sum-free, since $60 > 50 = \lceil
100/2 \rceil$ (kernel-certified).

The genuine boundary density is $\tfrac12$: sum-free sets exist up
to density exactly $\tfrac12$ (the odds) and none beyond.

\section{Verification}

\texttt{reproduce.py} checks the dense interval, the odd-number
sum-free set, and brute-forces the maximum sum-free size at
$N = 6, 8, 10$ ($= \lceil N/2 \rceil$).  The Lean package (core
Lean~4, v4.33.1) certifies \texttt{density\_below} ($24 < 25$),
\texttt{has\_sums}, \texttt{odds\_sumfree},
\texttt{ceiling\_exceeded} ($60 > 50$), and
\texttt{conjecture\_refuted}.  All five audited theorems report
\texttt{does not depend on any axioms}.

\section{Boundary}

The kernel certifies the density comparison, the sum witnesses, and
the ceiling violation; the $\lceil N/2 \rceil$ sum-free ceiling is
classical, cited in prose and brute-force confirmed.  The $5/8$
boundary claim is refuted under both readings; the true constant
is $1/2$.

\end{document}
